\documentclass{article}
\usepackage{amsmath}
\usepackage{amssymb}
\usepackage{fullpage}
\usepackage{enumerate}
\usepackage{xcolor}
\usepackage{bm}
\pagestyle{empty}
\usepackage{graphicx}
\graphicspath{{C:\Users\steve\Pictures}}
\begin{document}

\noindent{{\Large \bf Fall 2025, 4100 Problem Set 2, Hayden Fuller}

\large

\bigskip
\noindent{Please complete all problems. For any proof questions, please structure your proof similar to those from lecture.

\bigskip

\begin{enumerate}

\item Given $A=\left[
\begin{array}{cc}
3 & -1\\
0 & 2
\end{array}\right]$ show that the set $C=\{B\in\mathbb{R}^{2\times 2}:AB=BA\}$ is a subspaces of $\mathbb{R}^{2\times 2}$ and find a basis.
\\
\\assume $B=\left[\begin{array}{cc}a & b\\c & d\end{array}\right]$
\\$AB=\left[\begin{array}{cc}3 & -1\\0 & 2\end{array}\right]\left[\begin{array}{cc}a & b\\c & d\end{array}\right]=\left[\begin{array}{cc}3a-c & 3b-d\\2c & 2d\end{array}\right]$
\\$AB=\left[\begin{array}{cc}a & b\\c & d\end{array}\right]\left[\begin{array}{cc}3 & -1\\0 & 2\end{array}\right]=\left[\begin{array}{cc}3a & -a+2b\\3c & -c+2d\end{array}\right]$
\\assuming $AB=BA$, we have 
\\$3a-c=3a$, 
\\$3b-d=-a+2b$, 
\\$2c=3c$, 
\\$2d=-c+2d$
\\so $c=0$ and $d=a+b$
\\so $B=\left[\begin{array}{cc}a & b\\c & d\end{array}\right]=\left[\begin{array}{cc}a & b\\0 & a+b\end{array}\right]$
\\this shows that $C$ is a subspace since it contains zero and is closed under addition and scalar multiplication since $a$ and $b$ vary freely.
\\it is 2 dimensional, so a convenient basis is $\{\left[\begin{array}{cc}1 & 0\\0 & 1\end{array}\right],\left[\begin{array}{cc}0 & 1\\0 & 1\end{array}\right]\}$

\newpage\item Let $V$ be a vector space such that $dimV=n$. Show that any set of $n$ (non-zero) vectors, $\{{\bm v}_1,\dots,{\bm v}_n\}$, satisfying ${\bm v}_i^T{\bm v}_j=0$ for $i\neq j$ is a basis for $V$.
\\
\\any set of $n$ linearly independent vectors in an $n$-dimensional space is a basis
\\suppose $a_1v_1+a_2v_2+...+a_nv_n=0$
\\multiply both sides by $v_k$
\\$(a_1v_1+a_2v_2+...+a_nv_n=0)^Tv_k=a_1(v_1^Tv_k)+...+a_n(v_n^Tv_k)=0$
\\all terms are orthogonal and dissappear except $a_k(v_k^Tv_k)=0$
\\since $v_k$ is non zero, $a_k$ must be zero. This holds for each $k=1...n$, so all $a_k$ are zero
\\therefore all vectors are linearly independent, and since there are $n$ linearly independent vectors in $n$-dimensional space, they form a basis.

\newpage\item  Let $\mathbb{P}_n$ be the space of all polynomials over $\mathbb{R}$ of degree at most $n$. Use the Rank-Nullity theorem to show that the subset $S=\{g\in \mathbb{P}_n: f'+4f=g~\textrm{for some}~f\in \mathbb{P}_n\}$ is a subspace of $\mathbb{P}_n$ of dimension $n+1$.
\\
\\define function $T(f)=f'+4f$
\\$T$ is linear because $(af+bg)'+4(af+bg)=a(f'+4f)+b(g'+4g)$
\\the set $S=\{g\in\mathbb{P}_n:g=f'+4f for some f\in\mathbb{P}_n\}$ is the image of T, so $S=$im$T$
\\if $f\in$ker$T$ then $f'+4f=0$, so ker$T=\{0\}$
\\by rank nullity $\mathbb{P}_n=$dim ker$T+$dim im$T=(0)+$dim$S$
\\ so dim$S$=dim$\mathbb{P}_n=n+1$
\\thus $S$ is a subspace of $\mathbb{P}_n$ and dim$S=n+1$

\newpage\item Let $A\in\mathbb{R}^{m\times n}$ and $B\in\mathbb{R}^{n\times l}$ such that $AB{\bm x}=\textbf{0}$ for all ${\bm x}\in \mathbb{R}^l$. Prove that ${\rm rank}(A)+{\rm rank}(B)\leq n$ using the Rank-Nullity theorem.
\\
\\consider $A:\mathbb{R}^n\rightarrow \mathbb{R}^m$ and$B:\mathbb{R}^l\rightarrow\mathbb{R}^n$
\\$ABx=0$ for all $x\in\mathbb{R}^l$ means the composition $A\circ B$ is the zero map, so every vector in image of B is sent to zero by A
\\im$(B)\subseteq$ker$(A)$
\\rank$(B)=$dim im$(B) \leq$dim ker$(A)=$nullity$(A)$
\\$n=$rank$(A)+$nullity$(A)\geq$rank$(A)+$rank$(B)$
\\rank$(A)+$rank$(B)\leq n$


\newpage\item Let $V$ be a vector space with basis $\{{\bm v}_1,\cdots,{\bm v}_{n-1}\}$ and let $T\in\mathcal{L}(V,V)$ be defined by $T{\bm v}_i={\bm v}_{n-i}$. Show that $T^{-1}=T$ and that $T-I$ is not invetible. (Here $(T-I){\bm x}=T{\bm x}-{\bm x}$)
\\
\\$T^2v_i=T(T(v_i))=Tv_{n-i}=v_{n-(n-i)}=v_i$
\\since $T^2$ acts as the identity on every basis vector, $T^2=I$, so $T^{-1}=T$
\\
\\$(T-I)(v_i+v_{n-i})=Tv_i+Tv_{n-i}-(v_i+v_{n-i})=v_{n-i}+v_i-(v_i+v_{n-i})=0$
\\thus every nonzero  vector of form $v_i+v_{n-i}$ lies in ker$(T-I)$, so ker$(T-I)\neq0$
\\a linear map is invertable iff its kernel is $\{0\}$, so $T-I$ is not invertable

\newpage\item Consider the block-diagonal matrix $A=\left[
\begin{array}{cc}
B & 0\\
0 & C
\end{array}
\right]$. Show that ${\rm nullity}(A)={\rm nullity}(B)+{\rm nullity}(C)$.
\\
\\suppose $B\in\mathbb{R}^{m\times n}$ and $C\in\mathbb{R}^{p\times q}$
\\then $A$ is a block matrix of size $(m+p)\times(n+q)$
\\a vector in the domain of $A$ has the form $\left[\begin{array}{cc}x\\y\end{array}\right]$, $x\in\mathbb{R}^n$, $y\in\mathbb{R}^q$
\\$A\left[\begin{array}{cc}x\\y\end{array}\right]=\left[\begin{array}{cc}B & 0\\0 & C\end{array}\right]\left[\begin{array}{cc}x\\y\end{array}\right]=\left[\begin{array}{cc}Bx\\Cy\end{array}\right]$
\\so ker$(A)=\{\left[\begin{array}{cc}x\\y\end{array}\right]:Bx=0,Cy=0\}=$ker$(B)\oplus$ker$(C)$
\\dim ker$(A)=$dim ker$(B)+$dim ker$(C)$
\\nullity$(A)=$nullity$(B)+$nullity$(C)$

\end{enumerate}

\end{document}